\chapter{Theoretical background}

Through this study, we intend to explore the theoretical foundations of the hierarchical linear model, examine the principles and techniques of variational Bayesian inference, and maintain a self-contained structure that combines the necessary theoretical background with practical applications to real-world data. In doing so, we hope to contribute to a greater understanding and broader application of these methods in statistical modeling and data analysis.

\vspace{0.5cm}

This background will include fundamental concepts of linear regression, hierarchical modeling, Bayesian statistics, and variational inference. By presenting these concepts, we aim to equip readers with the essential foundation for understanding the subsequent analyses and discussions in this thesis.

\vspace{0.5cm}

\section{Bayesian Inference}

A statistical model is a set of probability distributions characterized by an unknown parameter vector $\bm{\theta}\in\bm{\Theta}$, which must be estimated using some statistical estimation technique. In this context, the core idea of Bayesian inference is to treat $\bm{\theta}$ as a random quantity, fully described by a probability distribution from which we can obtain its parameters, taking into account the information provided by the observations.

\vspace{0.5cm}

In this way, as stated by \cite{HOFF} Bayesian inference provides an estimation framework that yields parameter estimates with good statistical properties, a computational approach for obtaining those estimates, and flexibility in modeling. To formalize this, \cite{HOFF} defines a \textit{sample space} $\bm{\mathcal{Y}}$ as the set of all possible datasets from which a single dataset $\bm{y}$ is realized, and the \textit{parameter space} $\bm{\Theta}$ as the set of possible parameter values, from which we aim to identify the value that best represents the characteristics of the population.

\vspace{0.8cm}

The idea of Bayesian inference is based on assuming probabilistic and numerical beliefs about $\bm{y}$ and $\bm{\theta}$, expressed through probabilistic assumptions on $\bm{\mathcal{Y}}$ and $\bm{\Theta}$. For each value of $\bm{\theta}\in\bm{\Theta}$, a prior distribution $p(\bm{\theta})$ is specified, summarizing all prior beliefs held by the investigator. Additionally, for each pair $\bm{\theta}\in\bm{\Theta}$ and $\bm{y}\in\bm{\mathcal{Y}}$, the sampling model $p(\bm{y}\mid\bm{\theta})$ describes the assumed probabilistic behavior of the data-generating process.

\vspace{0.5cm}

Using \textit{Bayes's Theorem} we can update our knowledge about $\bm{\theta}$ based on the information provided by the dataset $\bm{y}$, as stated by \cite{GELMAN14}, as follows:

\begin{equation}
p(\bm{\theta}\mid\bm{y})=\frac{p(\bm{y}\mid\bm{\theta})p(\bm{\theta})}{\int_{\bm{\Theta}}p(\bm{y}\mid\bm{\theta})p(\bm{\theta})d\bm{\theta}}  
\end{equation}


in the continuous case, or: 

\begin{equation}
p(\bm{\theta}\mid \bm{y})=\frac{p(\bm{y}\mid\bm{\theta})p(\bm{\theta})}{\sum_{\bm{\theta}\in\bm{\Theta}}p(\bm{y}\mid\bm{\theta})p(\bm{\theta})d\bm{\theta}}
\end{equation}

for the discrete case. In either case, the posterior distribution $p(\bm{\theta}\mid\bm{y})$ can often be approximated as:

\begin{equation} 
\label{bayes1}
p(\bm{\theta}\mid\bm{y})\propto p(\bm{y}\mid\bm{\theta})p(\bm{\theta}).    
\end{equation}

\vspace{0.5cm}

In general, let $\bm{\theta}=(\theta_1,\cdots,\theta_p)$ denote the vector of $p$ unknown parameters and $\bm{y}=(y_1,\cdots,y_n)$ the dataset of $n$ observations. This methodology can also be used to derive the posterior predictive distribution (\cite{HOFF}), as follows:

\begin{equation} 
\label{bayes_pred}
p(\tilde{y}\mid\bm{y})=\mathsf{E}\left(p(\tilde{y}\mid\bm{y},\bm{\theta})\right)=\int_{\bm{\Theta}} p(\tilde{y}\mid\bm{y},\bm{\theta})p(\bm{\theta}\mid\bm{y}) d\bm{\theta}.  
\end{equation}


\section{Markov Chain Monte Carlo Methods (MCMC)}

As described in \cite{GAMERMAN}, in practice, point and interval estimates of $\bm{\theta}$, or any function of $\bm{\theta}$, denoted $g(\bm{\theta})$, are often obtained using MCMC sampling methods, this methods are based on drawing values from the posterior distribution $p(\bm{\theta}\mid\bm{y})$ and using Markov chains to approximate that distribution, as explained by \cite{GELMAN14} and \cite{HOFF}, let $\bm{\theta}$ the vector of unknown parameters of interest, and let $\bm{y}$ be the observed values from the sampling distribution $p(\bm{y}\mid\bm{\theta})$.


\newpage
Suppose we can generate $B$ numerical samples of $\bm{\theta}$ from the posterior distribution $p(\bm{\theta}\mid\bm{y})$.

\begin{equation}
\label{MCMC}
\bm{\theta}^{(1)},\cdots,\bm{\theta}^{(B)}\stackrel{\text{iid}}\sim p(\bm{\theta}\mid\bm{y})
\end{equation}


then the empirical distribution of the samples $\bm{\theta}^{(1)},\cdots,\bm{\theta}^{(B)}$ approximates the true posterior distribution as $B\rightarrow\infty$, by the law of large numbers, we can estimate any function $g(\bm{\theta})$ as:

\begin{equation} 
\label{MCMC_g}
\mathsf{E}(g(\bm{\theta)\mid\bm{y}})=\int_{\bm{\Theta}}g(\bm{\theta})p(\bm{\theta}\mid\bm{y})d\bm{\theta}\approx \frac{1}{B}\sum_{b=1}^{B}g(\bm{\theta}^{(b)}).
\end{equation}

\vspace{0.5cm}

This methodology also enables prediction of new observations and imputation of 
missing values through the posterior predictive distribution. In particular, 
inference is based on draws from \( p(\tilde{y}\mid \bm{y}) \), which can be 
approximated using Monte Carlo methods. Specifically, given posterior samples 
\( \bm{\theta}^{(1)}, \dots, \bm{\theta}^{(B)} \sim p(\bm{\theta}\mid \bm{y}) \), 
we generate replicated datasets 
\( \bm{\tilde{y}}^{(1)}, \dots, \bm{\tilde{y}}^{(B)} \), where each 
\( \bm{\tilde{y}}^{(b)} = (\tilde{y}_1^{(b)}, \dots, \tilde{y}_n^{(b)}) \) is drawn 
conditionally on the corresponding parameter value. For each iteration 
\( b = 1, \dots, B \), we sample
\begin{equation}
\bm{\tilde{y}}^{(b)} \stackrel{\text{iid}}{\sim} p(y \mid \bm{\theta}^{(b)}),
\end{equation}
and compute a summary statistic \( t\!\left(\bm{\tilde{y}}^{(b)}\right) \), such as 
the mean, median, or sample standard deviation.

\vspace{0.5cm}

From a probabilistic perspective, the posterior predictive distribution is given by
\begin{equation}
p(\tilde{y}\mid \bm{y})
= \int_{\bm{\Theta}} p(\tilde{y}\mid \bm{\theta}) \, p(\bm{\theta}\mid \bm{y}) 
\, d\bm{\theta}
= \mathsf{E}_{\bm{\theta}\mid \bm{y}}\!\left[p(\tilde{y}\mid \bm{\theta})\right],
\end{equation}
which can be approximated via Monte Carlo integration as
\begin{equation}
p(\tilde{y}\mid \bm{y})
\approx \frac{1}{B} \sum_{b=1}^{B} p(\tilde{y} \mid \bm{\theta}^{(b)}).
\end{equation}

\vspace{0.5cm}

In the same manner, posterior predictive distributions can be obtained for any 
statistic of interest. From a predictive perspective, this allows us to assess 
whether the observed statistic \( t(\bm{y}) \) is extreme relative to the 
posterior predictive distribution of \( t(\bm{\tilde{y}}) \). When 
\( t(\bm{y}) \) lies in the tail of 
\( p\!\left(t(\bm{\tilde{y}})\right) \), this indicates that the model may be 
unable to adequately reproduce important features of the observed data, thereby 
suggesting a potential lack of fit.




\subsubsection{Gibbs Sampling}

As described in \cite{HOFF}, Gibbs sampling is a widely used Markov Chain Monte Carlo (MCMC) algorithm for generating samples from a joint posterior distribution when the full conditional distributions (f.c.d.'s) of each parameter are available in closed form. Let 
$$
\bm{\theta} = \left( \theta_1, \ldots, \theta_p \right), \quad \bm{\theta}_{-j} = \left( \theta_1, \ldots, \theta_{j-1}, \theta_{j+1}, \ldots, \theta_p \right)
$$
be the vector of all parameters and the vector excluding the $j$th parameter, respectively. The full conditional distribution of $\theta_j$ is then:
\begin{equation} 
\label{fcd}
p\left(\theta_{j} \mid \bm{\theta}_{-j}, \bm{y}\right),
\end{equation}
from which we iteratively draw samples.

Given an initial value $\bm{\theta}^{(0)} = \left( \theta_1^{(0)}, \ldots, \theta_p^{(0)} \right)$, the Gibbs sampler produces a dependent sequence of vectors
$$
\bm{\theta}^{(1)}, \bm{\theta}^{(2)}, \ldots, \bm{\theta}^{(S)},
$$
where each $\bm{\theta}^{(s)}$ is obtained by successively sampling each coordinate from its corresponding f.c.d. conditioned on the most recent values of the other coordinates:
\begin{align*}
\theta_1^{(s)} &\sim p\left(\theta_1 \mid \theta_2^{(s-1)}, \theta_3^{(s-1)}, \ldots, \theta_p^{(s-1)}, \bm{y}\right), \\
\theta_2^{(s)} &\sim p\left(\theta_2 \mid \theta_1^{(s)}, \theta_3^{(s-1)}, \ldots, \theta_p^{(s-1)}, \bm{y}\right), \\
&\ \vdots \\
\theta_p^{(s)} &\sim p\left(\theta_p \mid \theta_1^{(s)}, \theta_2^{(s)}, \ldots, \theta_{p-1}^{(s)}, \bm{y}\right).
\end{align*}

This iterative process satisfies the \textit{Markov property}, since $\bm{\theta}^{(s)}$ depends on $\bm{\theta}^{(s-1)}$ only through the most recent updates of its coordinates, making $\left\{\bm{\theta}^{(s)}\right\}_{s \ge 0}$ a Markov chain whose stationary distribution is the target joint posterior. Under standard regularity conditions, the chain converges to the stationary distribution, enabling posterior inference.

\begin{algorithm}[H]
\caption{Gibbs Sampling}
\begin{algorithmic}
\State \textbf{Input:} Initial values $\theta_j^{(0)}$ for $j = 1,\dots,p$, data $\bm{y}$, number of iterations $B$
\For{$b = 1,\dots,B$}
  \For{$j = 1,\dots,p$}
    \State Sample $\theta_j^{(b)} \sim p\left(\theta_j \mid \theta_1^{(b)}, \ldots, \theta_{j-1}^{(b)}, \theta_{j+1}^{(b-1)}, \ldots, \theta_p^{(b-1)}, \bm{y}\right)$
  \EndFor
\EndFor
\State \textbf{Output:} Draws $\left\{\bm{\theta}^{(b)}\right\}_{b=1}^B$
\end{algorithmic}
\end{algorithm}



\subsubsection{Metropolis--Hastings}

When the full conditional distribution of a parameter does not have a closed form or is difficult to sample from directly, the Metropolis--Hastings (MH) algorithm provides a general approach for obtaining draws from the target distribution. A candidate value is generated from a proposal density and then accepted or rejected according to an acceptance ratio that corrects for differences between the proposal and the target posterior.

\begin{algorithm}[H]
\caption{Metropolis--Hastings (MH)}
\begin{algorithmic}
\State \textbf{Input:} initial values $\{\theta_j^{(0)}\}_{j=1}^p$, data $\bm{y}$, proposal density $q(\cdot\mid\cdot)$, number of iterations $B$
\For{$b = 1,\dots,B$}
  \For{$j = 1,\dots,p$}
    \State Generate a candidate value:
    \[
      \theta_j^{*(b)} \sim q\!\left(\theta_j^{*}\mid\theta_j^{(b-1)}\right)
    \]
    \State Compute the acceptance ratio:
    \[
      r(\theta_j) = 
      \frac{
        p\!\left(\theta_j^{*(b)} \mid \bm{\theta}_{-j}^{(b)},\bm{y}\right)\,
        q\!\left(\theta_j^{(b-1)}\mid\theta_j^{*(b)}\right)
      }{
        p\!\left(\theta_j^{(b-1)} \mid \bm{\theta}_{-j}^{(b)},\bm{y}\right)\,
        q\!\left(\theta_j^{*(b)}\mid\theta_j^{(b-1)}\right)
      }
    \]
    \State Draw $u \sim \mathsf{U}(0,1)$
    \If{$u < r(\theta_j)$}
      \State Accept: $\theta_j^{(b)} = \theta_j^{*(b)}$
    \Else
      \State Reject: $\theta_j^{(b)} = \theta_j^{(b-1)}$
    \EndIf
  \EndFor
\EndFor
\State \textbf{Output:} posterior draws $\{\bm{\theta}^{(b)}\}_{b=1}^B$
\end{algorithmic}
\end{algorithm}

\subsubsection{Hamiltonian Monte Carlo}

Hamiltonian Monte Carlo (HMC) is a Markov chain Monte Carlo method designed to overcome the limitations of random--walk behavior and slow convergence that often appear in classical approaches such as Gibbs sampling and Metropolis--Hastings \cite{GELMAN14}. By introducing an auxiliary momentum variable $\bm{\phi}$, HMC simulates Hamiltonian dynamics to generate proposals that explore the posterior distribution more efficiently. Specifically, samples are drawn from the joint distribution
\begin{equation}
p(\bm{\theta},\bm{\phi}\mid\bm{y}) = p(\bm{\phi})\,p(\bm{\theta}\mid\bm{y}),
\end{equation}
where $\bm{\phi}\sim \mathsf{N}_p(\bm{0},\bm{M})$ with mass matrix $\bm{M}$ of the same dimension as $\bm{\theta}$.


\begin{algorithm}[H]
\caption{Hamiltonian Monte Carlo (HMC)}
\begin{algorithmic}
\State \textbf{Input} Init values $\bm{\theta}^{(0)}$, step size $\epsilon$, number of leapfrog steps $L$, mass matrix $\bm{M}$, number of iterations $B$
\For{$b = 1,\dots,B$}
  \State Sample momentum $\bm{\phi}^{(b,0)} \sim \mathsf{N}_p(\bm{0},\bm{M})$
  \State Set $\bm{\theta}^{(b,0)}=\bm{\theta}^{(b-1)}$
  \For{$l = 1,\dots,L$}
    \State Update momentum $\bm{\phi}$
    $$
        \bm{\phi}^{(b,l-\frac{1}{2})}
        =
        \bm{\phi}^{(b,l-1)}+\frac{\epsilon}{2}\frac{\partial}{\partial\bm{\theta}}\log p\!\left(\bm{\theta}^{(b,l-1)}\mid \bm{y}\right)
        $$
    \State Update position $\bm{\theta}$ from eq. 
    $$
        \bm{\theta}^{(b,l)}
        =
        \bm{\theta}^{(b,l-1)}+\epsilon\bm{M}^{-1}\bm{\phi}^{(b,l-\frac{1}{2})}
        $$
    \State Update momentum $\bm{\phi}$
    $$
        \bm{\phi}^{(b,l)}
        =
        \bm{\phi}^{(b,l-\frac{1}{2})}+\frac{\epsilon}{2}\frac{\partial}{\partial\bm{\theta}}\log p\!\left(\bm{\theta}^{(b,l)}\mid \bm{y}\right)
        $$
  \EndFor
  \State Set $\bm{\theta}^{*}=\bm{\theta}^{(b,L)}$ and $\bm{\phi}^{*}=\bm{\phi}^{(b,L)}$
  \State Calculate acceptance ratio $r=\frac{p(\bm{\theta}^{*}\mid \bm{y})p(\bm{\phi}^{*})}{p(\bm{\theta}^{(b-1)}\mid \bm{y})p(\bm{\phi}^{(b,0)})}$ 
  \State Draw $u \sim \mathsf{U}(0,1)$
  \If{$u < \min\{1,r\}$}
    \State $\bm{\theta}^{(b)} = \bm{\theta}^*$
  \Else
    \State $\bm{\theta}^{(b)} = \bm{\theta}^{(b-1)}$
  \EndIf
\EndFor
\State \textbf{Output} Draws $\{\bm{\theta}^{(b)}\}_{b=1}^B$
\end{algorithmic}
\end{algorithm}


For further details on choosing $\epsilon$, $L$ and $\bm{M}$, see \cite{GELMAN14}.




\section{Bayesian Inference in Exponential Family of Distributions}

Consider $y_i\stackrel{\text{iid}}\sim p(y_i\mid\bm{\theta})$, where $p(\cdot)$ belongs to the exponential family of distributions. Thus

\begin{equation}
\label{dist_FED}
    p(y_i\mid\bm{\theta}) = h(y_i)\exp{\lbrace{\bm{g(\bm{\theta})}^\top \bm{t}\left(y_i\right)}-a(\bm{\theta})\rbrace},
\end{equation}

\newpage

hence, the likelihood function $p(\bm{y}\mid\bm{\theta})$ is given by

\begin{equation}
\label{loglik_FED}
p(\bm{y}\mid\bm{\theta}) =\underbrace{\left(\prod_{i=1}^{n}h(y_i)\right)}_{h(\bm{y})}\exp{\lbrace{\underbrace{\bm{g(\bm{\theta})}^\top}_{\bm{\eta}^T} \sum_{i=1}^{n}\bm{t}\left(y_i\right)}-na(\bm{\theta})\rbrace},
\end{equation}

reparameterizing \ref{loglik_FED} as a member of the natural exponential family \textsf{NEF} (\ref{NEF}), we obtain

\begin{equation}
\label{loglik_FEDN}
p(\bm{y}\mid\bm{\eta}) = h(\bm{y})\exp{\lbrace{\bm{\eta}}^\top \sum_{i=1}^{n}\bm{t}\left(y_i\right)-na(\bm{\eta})\rbrace}.
\end{equation}


\vspace{0.5cm}

Now, consider a prior distribution $p(\bm{\eta}\mid\bm{\lambda})$ for $\bm{\eta}$, also from the exponential family in its natural form:

\begin{equation}
\label{etaprior}
\begin{split}
    p(\bm{\eta}\mid\bm{\lambda}) =& h(\bm{\eta})\exp{\lbrace{\bm{\lambda}}^\top\bm{t}\left(\bm{\eta}\right)}-b(\bm{\lambda})\rbrace\\
    =&h(\bm{\eta})\exp{\lbrace{\bm{\eta}^{\top}\bm{\lambda}_{1}+\lambda_2(-na(\bm{\eta}))-b(\bm{\lambda})\rbrace}}
\end{split}
\end{equation}

where $\bm{\lambda}^\top= (\bm{\lambda}_{1}^\top,\lambda_2)$  and $\bm{t}(\bm{\eta})=(\bm{\eta},-na(\bm{\eta}))^\top$ with $\mathsf{dim}(\bm{t}(\bm{y}))=\mathsf{dim}(\bm{\lambda})=\mathsf{dim}(\bm{\eta})+1$. So the resulting posterior distribution is given by:

\begin{equation}
\label{etaposterior}
p(\bm{\eta}\mid\bm{y},\bm{\lambda})\propto\exp{\left\lbrace\left(\bm{\lambda}_{1}+\sum_{i=1}^{n}\bm{t}\left(y_i\right)\right)^\top\bm{\eta}+\left(\lambda_2+n\right)\left(-a(\bm{\eta})\right)\right\rbrace}
\end{equation}

it shows that the posterior distribution of $\bm{\eta}$ remains in the same exponential family as the prior, due to the conjugacy. Therefore, the natural parameter of the posterior distribution, $\bm{h}^{\top}=(\bm{h}_{1}^{\top},h_{2})$ is given by:

\begin{equation}
    \bm{h}_{1}=\bm{\lambda}_1+\sum_{i=1}^{n}\bm{t}\left(y_i\right)\quad\text{and}\quad h_{2}=\lambda_2+n
\end{equation}

this result can be extended to derive the natural parameters of posterior distributions in Bayesian hierarchical models. For more details, see \cite{LAWRENCE}.

\section{Bayesian Variational Inference}

In \cite{ZETER} it mentions that the variational approach has a different main idea than MCMC methods, instead of get samples from a distribution, the variational approach try to determine a probability distribution $q(\bm{\theta})\in\bm{\mathcal{Q}}$ with $\bm{\mathcal{Q}}=\left\{q(\bm{\theta},\bm{\lambda}): \bm{\lambda}\in\bm{\Lambda}\right\}$ where $\bm{\lambda}$ are the parameters of variational distribution,  which is as close as possible to our probability distribution of interest $p(\bm{\theta}\mid\bm{y})$, which gives as result an optimization problem that needs specify a cost function and similarity metric to measure the closeness between these two probability distributions.
For this purpose, we will choose $q^{*}(\bm{\theta})$ such that:

\begin{equation}
\label{argmin_vi}
q^{*}(\bm{\theta})=
\underset{q(\bm{\theta}) \in \bm{\mathcal{Q}}}{\arg\min}
\lbrace\bm{D}\left(
q(\bm{\theta}),
p(\bm{\theta}\mid\bm{y})
\right)\rbrace
\quad
\Leftrightarrow
\quad
\bm{\lambda}^{*}
=
\underset{\bm{\lambda}\in\bm{\Lambda}}{\arg\min}
\lbrace\bm{D}\left(
q(\bm{\theta},\bm{\lambda}),
p(\bm{\theta}\mid\bm{y})
\right)\rbrace,
\end{equation}

with $\bm{D}(\cdot,\cdot)$ the similarity metric to be minimized, in this case the Bayesian Variational inference purposes use the Kullback-Leibler divergence \ref{KL_div} as metric to measure the closeness between probability distirbutions, as follows

\begin{equation}
    \label{VB}
    \mathsf{KL}\left(q(\bm{\theta})\mid\mid\ p(\bm{\theta}\mid\bm{y})\right)=\int q(\bm{\theta}) \log \left(\frac{q(\bm{\theta})}{p(\bm{\theta}\mid\bm{y})}\right)d\bm{\theta}
\end{equation}

\subsection{Kullback-Leibler Divergence}

The Kullback-Leibler (KL) divergence \cite{KL} is a statistical measurement from information theory that is commonly used to quantify the difference between a probability distribution $q(\bm{y})$ and the reference distribution $p(\bm{y})$. It is defined as:

\begin{equation}
\label{KL_div}
\mathsf{KL}\left(p(\bm{y})\mid\mid q(\bm{y})\right)=\mathsf{E}_{p(\bm{y})}\left[\log\left(\frac{p(\bm{y})}{q(\bm{y})}\right)\right]
\end{equation}

that is, 

\begin{equation}
\label{KL_div1}
\mathsf{KL}\left(p(\bm{y})\mid\mid q(\bm{y})\right)=\int_{-\infty}^{\infty}p(\bm{y})\log\left(\frac{p(\bm{y})}{q(\bm{y})}\right)d\bm{y}
\end{equation}

in the continuous case, or in the discrete case:

\begin{equation}
\label{KL_div2}
\mathsf{KL}\left(p(\bm{y})\mid\mid q(\bm{y})\right)=\sum_{y \in \bm{y}} p(\bm{y}) \log\left(\frac{p(\bm{y})}{q(\bm{y})}\right)
\end{equation}

\subsubsection{Properties}

\begin{itemize}
    \item $\mathsf{KL}\left(p(\bm{y})\mid\mid q(\bm{y})\right)\geq0$
    \item $\mathsf{KL}\left(p(\bm{y})\mid\mid q(\bm{y})\right)=0$ if and only if $p(\bm{y})=q(\bm{y})$
    \item $\mathsf{KL}\left(p(\bm{y})\mid\mid q(\bm{y})\right)\neq \mathsf{KL}\left(q(\bm{y})\mid\mid p(\bm{y})\right)$
\end{itemize}

\subsection{Evidence Lower Bound (ELBO)}

As described in \cite{Ormerod2010-bc}, consider a Bayesian model with $\bm{\theta}\in\bm{\Theta}$ and observed data $\bm{y}\in\bm{\mathcal{Y}}$. The posterior distribution for $\bm{\theta}$ is given by:

\begin{equation}
\label{post_theta}
p(\bm{\theta}\mid\bm{y})=\frac{p(\bm{y},\bm{\theta})}{p(\bm{y})}
\end{equation}

where $p(\bm{y})$ is the marginal likelihood or model evidence. Introducing an arbitrary probability density function $q(\bm{\Theta})$ over $\bm{\Theta}$, we can rewrite the marginal likelihood as:

\begin{equation}
\label{lnPy}
\begin{split}
    \log({p(\bm{y})})=& \log (p(\bm{y}))\int q(\bm{\theta}) d\bm{\theta}\\
                    =& \int q(\bm{\theta}) \log (p(\bm{y})) d\bm{\theta}\\
                    =& \int q(\bm{\theta}) \log \left(\frac{p(\bm{y},\bm{\theta})}{p(\bm{\theta}\mid\bm{y})}\right) d\bm{\theta}\\
                    =& \int q(\bm{\theta}) \log \left(\frac{p(\bm{y},\bm{\theta})/q(\bm{\theta})}{p(\bm{\theta}\mid\bm{y})/q(\bm{\theta})}\right) d\bm{\theta}\\
                    =&  \int q(\bm{\theta}) \log \left(\frac{p(\bm{y},\bm{\theta})}{q(\bm{\theta})}\right)d\bm{\theta}+ \int q(\bm{\theta}) \log \left(\frac{q(\bm{\theta})}{p(\bm{\theta}\mid\bm{y})}\right)d\bm{\theta}\\
                    =&\int q(\bm{\theta}) \log \left(\frac{p(\bm{y},\bm{\theta})}{q(\bm{\theta})}\right)d\bm{\theta}+\underbrace{\mathsf{KL}\left(q(\bm{\theta})\mid\mid\ p(\bm{\theta}\mid\bm{y})\right)}_{\geq 0}\\
                    \geq&\int q(\bm{\theta}) \log \left(\frac{p(\bm{y},\bm{\theta})}{q(\bm{\theta})}\right)d\bm{\theta}\\
                    \geq&\quad \mathsf{E}_{q(\bm{\theta})}\left[ \log \left(\frac{p(\bm{y},\bm{\theta})}{q(\bm{\theta})}\right)\right]\\
                    \geq&\quad \mathcal{L}(q)
\end{split}
\end{equation}

from the final inequality, it follows that: 

\begin{equation}
\label{ELBO}
\begin{split}
    \mathcal{L}(q) =&\mathsf{E}_{q(\bm{\theta})}\left[ \log \left(\frac{p(\bm{y},\bm{\theta})}{q(\bm{\theta})}\right)\right]\\
    =&\mathsf{E}_{q(\bm{\theta})}\left[\log (p(\bm{y},\bm{\theta})) \right]-\mathsf{E}_{q(\bm{\theta})}\left[\log (q(\bm{\theta}) )\right]\\
    =&\int q(\bm{\theta}) \log \left(\frac{p(\bm{y},\bm{\theta})}{q(\bm{\theta})}\right)d\bm{\theta}\\
    =&-\mathsf{KL}\left(q(\bm{\theta})\mid\mid p(\bm{y},\bm{\theta})\right)
\end{split}
\end{equation}

this quantity $\mathcal{L}(q)$ is known as \textbf{\textit{Evidence Lower Bound (ELBO)}} and serves as a lower bound for $\log{(p(\bm{y}))}$, note that $\log({p(\bm{y})})$ does not depend on $q(\bm{\theta})$ and then the problem of minimize \ref{VB} is equivalent to maximize $\mathsf{ELBO}$ keeping in mind that $\log({p(\bm{y})})\geq \mathsf{E}_{q(\bm{\theta})}\left[ \log \left(\frac{p(\bm{y},\bm{\theta})}{q(\bm{\theta})}\right)\right]$ as in \ref{lnPy}.


\subsection{Mean Field Approach}

The mean field approach assumes that the elements of $\bm{\theta}$ are independent, introducing a factorization of the variational distribution $q(\bm{\theta})$ as follows:

\begin{equation}
\label{MFA}
\begin{split}
       q(\bm{\theta})
       =&\prod_{j=1}^{p}q\left(\theta_j\right) 
\end{split}
\end{equation}

this assumption allows for a tractable solution for each $\theta_j$ and facilitates efficient optimization. The approach is effective in capturing marginal effects; however, it tends to underestimate the variance and higher-order moments, as it does not account for correlations among the parameters.

\subsection {Coordinate Ascent Variational Inference}

To optimize the  $\mathsf{ELBO}$ in Equation \ref{ELBO}, the  Coordinate Ascent Variational Inference ($\mathsf{CAVI}$) algorithm iteratively optimizes each factor of the variational distribution in equation \ref{MFA}, holding the other components fixed. This process resembles Gibbs Sampling, utilizing the chain rule of probability to write the joint distribution $p(\bm{y},\bm{\theta})$ as:

\begin{equation}
    p(\bm{y},\bm{\theta}) = p(\bm{y})\prod_{j=1}^{p}p(\theta_j\mid \bm{\theta}_{-j},\bm{y})
\end{equation}

Consider $\mathcal{L}(q)$ as a function of $q(\theta_j)=q(\theta_j,\bm{\lambda}_j)$ for some $j\in\lbrace1,\cdots,p\rbrace$, while keeping the other $p-1$ factors fixed. Then equation \ref{ELBO} can be rewritten as:

\begin{equation}
\label{CAVI}
\begin{split}    
    \mathcal{L}(q)=&\mathsf{E}_{q(\bm{\theta})}\left[\log (p(\bm{y},\bm{\theta})) \right]-\mathsf{E}_{q(\bm{\theta})}\left[\log (q(\bm{\theta}) )\right]\\
                  =&\log \left(p(\bm{y})\right) + \sum_{j=1}^{p}\mathsf{E}_{q(\bm{\theta})}\left[\log\left(p(\theta_j\mid \bm{\theta}_{-j},\bm{y})\right)\right]-\sum_{j=1}^{p}\mathsf{E}_{q(\bm{\theta})}\left[\log\left(q(\theta_j)\right)\right]\\
                  =& \mathsf{E}_{q(\bm{\theta})}\underbrace{\left[\log p(\theta_j\mid \bm{\theta}_{-j},\bm{y})\right]}_{\log{p(\bm{y},\bm{\theta})}-\log{p(\bm{\theta}_{-j},\bm{y})}}-\mathsf{E}_{q(\bm{\theta})}\left[\log\left(q(\theta_j)\right)\right] + \sum_{i\neq j}^{p}\left[\mathsf{E}_{q(\bm{\theta})}\left[\log\left(p(\theta_i\mid \bm{\theta}_{-i},\bm{y})\right)\right]-\mathsf{E}_{q(\bm{\theta})}\left[\log\left(q(\theta_i)\right)\right]\right] +\text{const} \\
                  =&\int q\left(\theta_j\right)\underbrace{\bigl\{\int \log{p(\bm{y},\bm{\theta})}\prod_{i\neq j}q\left(\theta_i\right) d\bm{\theta}_{-j}\bigr\}}_{{\mathsf{E}_{-j}}\left[\log{p(\bm{y},\bm{\theta})}\right]}d\theta_j-\int q\left(\theta_j\right) \log q\left(\theta_j\right)\bigl\{\int \prod_{i\neq j}q\left(\theta_i\right) d\bm{\theta}_{-j}\bigr\}d\theta_j+\text{const}\\
                  =&\int q\left(\theta_j\right)\bigl\{\mathsf{E}_{-j}\left[\log{p(\bm{y},\bm{\theta})}\right]-\log q\left(\theta_j\right)\bigr\}d\theta_j+\text{const}\\
                  =& -\mathsf{KL}\left(q\left(\theta_j\right)\mid\mid\exp{\lbrace\mathsf{E}_{- j}\left[\log{p(\bm{y},\bm{\theta})}\right]}\rbrace\right)+\text{const}\\
\end{split}
\end{equation}

where $\mathsf{E}_{j}\left[\cdot\right]$ and $\mathsf{E}_{-j}\left[\cdot\right]$ denote expectations taken with respect to the variational distributions $q(\theta_j)$ and $q(\bm{\theta}_{-j})$, respectively. Specifically, $\mathsf{E}_{j}\left[\cdot\right]$ represents expectation over the variational distribution associated with the $j$-th model parameter $\theta_j$, while $\mathsf{E}_{-j}\left[\cdot\right]$ corresponds to expectation with respect to the joint variational distribution of all remaining parameters $\bm{\theta}_{-j} = \{\theta_1,\dots,\theta_{j-1},\theta_{j+1},\dots,\theta_p\}$.

\vspace{0.5cm}

In the context of Coordinate Ascent Variational Inference, the expectation $\mathsf{E}_{-j}\left[\log p(\bm{y},\bm{\theta})\right]$ plays a central role, it represents the expected log-joint distribution of the data and parameters, averaged over the current variational approximations of all parameters except $\theta_j$. This operation effectively isolates the contribution of $\theta_j$ to the joint density while holding the remaining variational factors fixed.

\newpage

The final expression in Equation~\ref{CAVI} can therefore be interpreted as the minimization of a Kullback-Leibler divergence between the variational factor $q(\theta_j)$ and an unnormalized target distribution attains its minimum when
$$
\mathsf{KL}\left(
q(\theta_j)\;\middle\|\;
\exp\left\{\mathsf{E}_{-j}\left[\log p(\bm{y},\bm{\theta})\right]\right\}
\right)=0,
$$

which yields the optimal functional form of the variational factor $q(\theta_j)$ within the chosen variational family.


\begin{equation}
\begin{split}
\label{MVFB_opt}
   \log q\left(\theta_j\right)=& \mathsf{E}_{{- j}}\left[\log{p(\bm{y},\bm{\theta})}\right]\\
             q\left(\theta_j\right)=&\frac{\exp{\left(\mathsf{E}_{{-j}}\left[\log{p(\bm{y},\bm{\theta})}\right]\right)}}{\int \exp{\left(\mathsf{E}_{{- j}}\left[\log{p(\bm{y},\bm{\theta})}\right]\right)} d\theta_j}\\
            q\left(\theta_j\right)\propto& \exp{\left(\mathsf{E}_{{- j}}\left[\log{p(\bm{y},\bm{\theta})}\right]\right)}
\end{split}    
\end{equation}

this expression corresponds to the optimal form under the mean-field approximation \cite{BISHOP}. Each variational distribution associate with each $\theta_j$ model parameters depends on the expectations under the remaining \( p-1 \) variational distributions and expected values of its parameters, resulting in a set of coupled equations. These interdependencies define the structure of the iterative Coordinate Ascent Variational Inference (CAVI) (algorithm \ref{CAVI_alg}).

\begin{algorithm}[H]
\caption{Coordinate Ascent Variational Inference (CAVI)}
\label{CAVI_alg}
\begin{algorithmic}
\Statex \textbf{Input:} Set $i = 0$, initial variational parameters 
$\bm{\lambda}_j^{(0)}$ for $j = 1,\dots,p$, observed data $\bm{y}$

\While{ELBO has not converged}
    \For{$j = 1,\dots,p$}
        \State  Update $q\left(\theta_j,\bm{\lambda}_j^{(i+1)}\right)\propto\exp{\left(\mathsf{E}_{{-j}}^{(i)}\left[\log{p(\bm{y},\bm{\theta})}\right]\right)}$ and extract its parameters
        $\bm{\lambda}_j^{(i+1)}$
    \EndFor
    \State Compute the ELBO:
    $$
    \mathcal{L}^{(i+1)}\left(q\right)
    =
    \mathsf{E}^{(i+1)}_{q(\bm{\theta})}\left[\log p(\bm{y},\bm{\theta})\right]
    -
    \mathsf{E}^{(i+1)}_{q(\bm{\theta})}\left[\log q(\bm{\theta})\right]
    $$
    \State Set $i =  i + 1$
\EndWhile
\Statex \textbf{Output:} Optimal variational parameters 
$\bm{\lambda}_j$ for $j = 1,\dots,p$
\end{algorithmic}
\end{algorithm}

To ensure tractability and computational efficiency, it is customary to restrict the variational distribution to a family $q(\bm{\theta}) \in \bm{\mathcal{Q}}$, where $\bm{\mathcal{Q}}$ denotes a class of tractable probability distributions, as discussed in \cite{HOFFMAN}. Under this restriction, the variational inference problem reduces to maximizing the $\mathsf{ELBO}$ over the chosen family. In particular, the variational distribution $q(\bm{\theta})$ is parameterized as a member of NEF ~\eqref{NEF}, that is, $q(\bm{\theta}) = q(\bm{\theta}, \bm{\nu})$, where $\bm{\nu}$ denotes the vector of natural parameters associated with the variational family.

\newpage

Notice that $p(\bm{y},\bm{\theta})  \propto p(\theta_j, \bm{\theta}_{-j},\bm{y}) \propto p(\theta_j\mid \bm{\theta}_{-j},\bm{y})$, and the denominator of $p(\theta_j\mid \bm{\theta}_{-j},\bm{y})$ does not depend on $\theta_j$, the optimal solution \ref{MVFB_opt} becomes:

\begin{equation}
\label{MVFB_opt_post}
q\left(\theta_j\right)\propto \exp{\left(\mathsf{E}_{{- j}}\left[\log{p(\theta_j\mid \bm{\theta}_{-j},\bm{y})}\right]\right)}
\end{equation}

if $p(\theta_j\mid \bm{\theta}_{-j},\bm{y})$ belongs to the \textsf{NEF} (\ref{NEF}), then:

\begin{equation}
\label{fullconditionalfed}
    p(\theta_j\mid \bm{\theta}_{-j},\bm{y}) = h(\theta_j)\exp{\lbrace\bm{\eta}_j( \bm{\theta}_{-j},\bm{y})^\top\bm{t}(\theta_{j})-a\left(\bm{\eta}_j( \bm{\theta}_{-j},\bm{y})\right)\rbrace}
\end{equation}

substituting into \ref{MVFB_opt},

\begin{equation}
   \mathsf{E}_{{-j}}\left[ \log p(\theta_j\mid \bm{\theta}_{-j},\bm{y})\right] =\mathsf{E}_{{-j}}\left[ \log h(\theta_j)\right]+\mathsf{E}_{{-j}}\left[\bm{\eta}_j( \bm{\theta}_{-j},\bm{y})^\top\bm{t}(\theta_{j})-a\left(\bm{\eta}_j( \bm{\theta}_{-j},\bm{y})\right)\right]
\end{equation}

thus, the optimal update becomes:

\begin{equation}
   q\left(\theta_j\right) \propto h(\theta_j)\exp{\lbrace\mathsf{E}_{{-j}}\left[\bm{\eta}_j( \bm{\theta}_{-j},\bm{y})\right]^\top\bm{t}(\theta_{j})\rbrace}
\end{equation}

which implies that each $q(\theta_j)=q(\theta_j,\nu_j)$ also belongs to the \textsf{NEF} (\ref{NEF}), with natural parameter: 

\begin{equation}
\label{NMVFB_opt}
\nu_{j}=\mathsf{E}_{{-j}}\left[\bm{\eta}_j( \bm{\theta}_{-j},\bm{y})\right]
\end{equation}

this result shows the direct relationship between the variational distribution and its full conditional distribution, providing a straightforward mechanism for updating natural parameters when the full conditional distribution is known.

\subsection{Stochastic Variational Inference}
\label{SVI}
One of most important type of exponential family models are the \textbf{\textit{conditionally conjugate models}} \cite{Blei2016VariationalIA}, where the parameter vector $\bm{\theta}$ is partitioned into two sets: $\bm{\theta}=(\bm{\theta_g},\bm{\theta_l})$. Here, $\bm{\theta_g}$ represents global latent variables, associated with the overall model structure, while $\bm{\theta_l}$ corresponds to local latent variables, one per data point \cite{REGUEIRO}. In this framework, the joint distribution factorizes as

\begin{equation}
    p(\bm{y},\bm{\theta}) = p(\bm{\theta_g},\bm{\theta_l},\bm{y})   = p(\bm{\theta_{l}},\bm{y}\mid \bm{\theta_g})p(\bm{\theta_g}) = p(\bm{\theta_g})\prod_{i=1}^{n}p(\theta_{l_i},y_i\mid\bm{\theta_g})
\end{equation}

where both $p(\bm{\theta_{l}},\bm{y}\mid \bm{\theta_g})$ and $p(\bm{\theta_g})$ are conjugate distributions in the exponential family. Specifically

\begin{equation}
    p(\theta_{l_i},y_i\mid\bm{\theta}_g)\propto\exp{\left\lbrace\bm{\theta}_g^\top\bm{t}(\theta_{l_i},y_i)-a_l(\bm{\eta})\right\rbrace},
\end{equation}

with prior distribution for $\bm{\theta}_g$ given by

\begin{equation}
    p(\bm{\theta}_g)\propto\exp{\left\lbrace\bm{\lambda}^\top\bm{t}(\bm{\theta}_g)-a_g(\bm{\lambda})\right\rbrace},
\end{equation}

here $\bm{\lambda}=(\bm{\lambda}_{1},\lambda_{2})^\top$ is the natural parameter. By conjugacy, $\bm{t}(\bm{\theta}_g) = (\bm{\theta}_g,-a_l(\bm{\theta}_g))^\top$ which implies that the full conditional distribution of $\bm{\theta}_g$ belongs to same exponential family as the prior,

\begin{equation}
    p(\bm{\theta}_g \mid \bm{\theta_l},\bm{y} ) \propto \exp{\left\lbrace\bm{\eta}_g(\bm{\theta_l},\bm{y})^\top\bm{t}(\bm{\theta}_g)\right\rbrace},
\end{equation}

where the natural parameter is

\begin{equation}
    \bm{\eta}_g(\bm{\theta_l},\bm{y})=\bm{\lambda}+\sum_{i=1}^{n}(\bm{t}\left(\theta_{l_i},y_i\right),1)^\top.
\end{equation}

\vspace{0.5cm}

Moreover the variational distribution is given by

\begin{equation}
    q(\bm{\theta})=q(\bm{\theta}_l)q(\bm{\theta}_g),
\end{equation}

so, applying the result \ref{NMVFB_opt}, the variational distribution for $\bm{\theta}_g$ parametrized in its NEF form (\ref{NEF}), becomes in

\begin{equation}
q(\bm{\theta}_{g}\mid\bm{h})\propto\exp{(\underbrace{\mathsf{E}_{l}\left[\bm{\eta}_g(\bm{\theta}_l,\bm{y})\right]^\top}_{\bm{h}^\top}\bm{\theta}_{g})},
\end{equation}

thus, the \textbf{global variational parameters} are given by

\begin{equation}
\label{global_updates}
\bm{h}=\bm{\lambda}+\sum_{i=1}^{n}(\mathsf{E}_{l}\left[\bm{t}\left(\theta_{l_i},y_i\right)\right],1)^\top.
\end{equation}

\vspace{0.5cm}

Now, assume that the full conditional for each $\theta_{l_i}$ also belongs to the NEF (\ref{NEF}), so,
\begin{equation}
        p(\theta_{l_i}\mid y_i,\bm{\theta}_g)\propto\exp{\left\lbrace{\bm{\eta}_l(\bm{\theta}_g,y_i)}^\top\theta_{l_i}-a\left(\bm{\eta}_l(\bm{\theta}_g,y_i)\right)\right\rbrace},
\end{equation}

then, the corresponding local variational distribution is

\begin{equation}
q(\theta_{l_i}\mid\bm{\varphi_i})\propto\exp{(\underbrace{\mathsf{E}_{g}\left[\bm{\eta}_l(\bm{\theta}_g,y_i)\right]^\top}_{\bm{\varphi_i}^\top}\theta_{l_i})},
\end{equation}

so the \textbf{local variational parameters} are given by

\begin{equation}
    \label{local_updates}
\bm{\varphi_i}=\mathsf{E}_{g}\left[\bm{\eta}_l(\bm{\theta_g},y_i)\right],
\end{equation}

note that the global variational parameters \ref{global_updates} depend on the expectations over local variational parameters \ref{local_updates}, creating a bidirectional dependency that drives the \textsf{CAVI} algorithm (Algorithm \ref{CAVI_FEDN_alg}).

\begin{algorithm}[H]
\caption{CAVI with conditionally conjugate models}
\begin{algorithmic}
\Statex \textbf{Input} Init values $\bm{h}^{(0)},\bm{\varphi}^{(0)}$ , data values $\bm{y}$, set $i=0$ 
\While{ELBO has not convergency}
    \For{$i = 1,...,n$}
        \State Optimize the local variational parameters $\bm{\varphi_i}^{(i+1)}=\mathsf{E}_{g}^{(i)}\left[\bm{\eta}_l(\bm{\theta_g},y_i)\right]$
    \EndFor
    \State Update the global variational parameters $\bm{h}^{(i+1)}=\bm{\lambda}+\sum_{i=1}^{n}(\mathsf{E}_{l}^{(i+1)}\left[\bm{t}\left(\theta_{l_i},y_i\right)\right],1)^\top$
    \State Calculate ELBO $
    \mathcal{L}^{(i+1)}\left(q\right)
    =
    \mathsf{E}^{(i+1)}_{q(\bm{\theta})}\left[\log p(\bm{y},\bm{\theta})\right]
    -
    \mathsf{E}^{(i+1)}_{q(\bm{\theta})}\left[\log q(\bm{\theta})\right]
    $
    \State Set $i=i+1$.
\EndWhile
\Statex \textbf{Output} Optimal variational natural parameters $\bm{h},\bm{\varphi}$
\end{algorithmic}
\label{CAVI_FEDN_alg}
\end{algorithm}

Despite its flexibility, the \textsf{CAVI} algorithm suffers from a significant limitation: each iteration requires a full pass through all data points, which limits scalability for large datasets. To overcome this, \cite{HOFFMAN} proposed \textit{Stochastic Variational Inference (SVI)}. 

\vspace{0.5cm}

\subsubsection {Stochastic Optimization}

\textit{SVI} optimizes the $\mathsf{ELBO}$ using stochastic gradient methods, relying on an unbiased estimator of the \textit{natural gradient} of the $\mathsf{ELBO}$. Instead of computing full datasets expectations at every step, \textit{SVI} uses mini-batches or single data points to estimate the gradient, combining classical  \textsf{CAVI} updates with the stochastic approximation rule from \cite{Robbins&Monro:1951}. This method ensures convergence under mild conditions and greatly improves computational efficiency, enabling variational inference on massive datasets.

\vspace{0.5cm}

The main idea of stochastic optimization is, given a function $f(x)$ to be maximized and $h_t(x)$ a realization of a random variable $H(x)$ whose expected value is $f(x)$, i.e., $\mathsf{E}\left[H(x)\right]=\nabla f(x)$ (\cite{RANGAN}) and then setting $\rho_t$ a nonnegative scalar, in such a way that the value $x_t$ which maximizes $f(x)$ is updated as follows

\begin{equation}
    \label{stoch_opt}
    x_t=x_{t-1}+\rho_th_t(x_t),
\end{equation}

this converges to a maximum of $f(x)$ when $\rho_{t}$ known as learning rate satisfies the Robbins-Monro  conditions,

\begin{equation*}
    \sum_{t=1}^{\infty}\rho_t=\infty,\quad
    \sum_{t=1}^{\infty}\rho_t^2<\infty,
\end{equation*}

\vspace{0.5cm}

a commonly used learning rate schedule, as suggested by \cite{HOFFMAN}, is

\begin{equation}
\rho_t=(t+\tau)^{-\chi} ,   
\end{equation}

where $\tau\geq0$ is a delay parameter, which downweights early iterations to improve stability, and $\chi\in\left(0.5,1\right]$ is the forgetting rate, which controls how quickly past updates are discounted.

\subsubsection{Natural Gradients of ELBO}

In this context, the natural gradient of the ELBO \ref{ELBO} with respect to the global variational parameters $ \bm{h} $ is given by:

\begin{equation}
\widetilde{\nabla}_{\bm{h}} \mathcal{L}(q) = \mathsf{E}_{l}\left[\bm{\eta}_g(\bm{\theta}_{g}, \bm{y})\right] - \bm{h}
\end{equation}

similarly, for each local variational parameter $ \bm{\varphi}_n $, the natural gradient is given by:

\begin{equation}
\widetilde{\nabla}_{\bm{\varphi}_i} \mathcal{L}(q) = \mathsf{E}_{g}\left[\bm{\eta}_l(\bm{\theta}_l, y_i)\right] - \bm{\varphi}_i
\end{equation}

These gradients represent the direction of steepest ascent in the Riemannian geometry induced by the exponential family and are often more efficient than standard gradients in variational inference optimization. For more theoretical background on natural gradients, see \cite{AMARI}.

\subsubsection{Noisy Gradient Estimator}

In practice, computing the exact gradient of the ELBO over the entire dataset is often computationally prohibitive. To address this, \textit{SVI} uses a mini-batch approach. Specifically, a random subsample $ \mathcal{S} \subset \{1, \dots, \mathcal{N}\} $ of size $ |\mathcal{S}|$ is selected, where $\mathcal{N}$ is the total number of sampling units (e.g. observations, groups) and elements $\mathcal{S}$ are sampled uniformly at random. 

\vspace{0.5cm}

To ensure that the gradient estimator remains unbiased with respect to the full dataset, a scaling factor is introduced. This factor adjusts the contribution of the mini-batch to reflect the entire dataset:

\begin{equation}
\mathcal{C} = \frac{\mathcal{N}}{|\mathcal{S}|}
\end{equation}

using this factor, the unbiased estimator of the natural gradient  with respect to $ \bm{h} $, known as \textbf{noisy gradient} is given by:

\begin{equation}
\widehat{\widetilde{\nabla}}_{\bm{h}} \mathcal{L}(q) =\bm{\lambda}+ \mathcal{C} \sum_{i \in \mathcal{S}}(\mathsf{E}_{l}\left[\bm{t}\left(\theta_{l_i},y_i\right)\right],1)^\top - \bm{h}
\end{equation}

so,

\begin{equation}
\mathsf{E}_{\mathcal{S}} \left[ \widehat{\widetilde{\nabla}}_{\bm{h}} \mathcal{L}(q) \right] = \widetilde{\nabla}_{\bm{h}} \mathcal{L}(q)
\end{equation}

this stochastic estimation enables scalable updates of the global variational parameters without the need to process the full dataset at each iteration, making \textit{SVI} particularly suitable for large-scale Bayesian inference problems.

\subsubsection{Robbins-Monro Update Rule}

Let $ \bm{h}^{(t)} $ denote the value of the global variational parameter at iteration $ t $, and let $ \rho_t \in (0,1) $ be the learning rate. The update rule for $ \bm{h} $ follows the stochastic approximation scheme:

\begin{equation}
\label{Robbins-Monro}
\begin{split}
\bm{h}^{(t)} =& \bm{h}^{(t-1)} + \rho_t \widehat{\widetilde{\nabla}}_{\bm{h}}\mathcal{L}(q(\bm\theta^{(t)}))\\
=& \bm{h}^{(t-1)} + \rho_t \left(\bm{\lambda}+\mathcal{C}\sum_{i \in \mathcal{S}}\left(\mathsf{E}_{l}\left[\bm{t}\left(\theta_{l_i}^{(t)},y_i\right)\right],1\right)^\top-\bm{h}^{(t-1)}\right)\\
 =& (1-\rho_t) \bm{h}^{(t-1)} + \rho_t \left(\bm{\lambda}+\mathcal{C}\sum_{i \in \mathcal{S}}\left(\mathsf{E}_{l}\left[\bm{t}\left(\theta_{l_i}^{(t)},y_i\right)\right],1\right)^\top\right) 
\end{split}
\end{equation}

\begin{algorithm}[H]
\caption{Stochastic Variational Inference (SVI)}
\begin{algorithmic}
\Statex \textbf{Input} Init values $\bm{h}^{(0)},\bm{\varphi}^{(0)}$, data values $\bm{y}$, learning rate parameters $\tau,\chi$, minibatch size $|\mathcal{S}|$, number of iterations $T$
\For{$t = 1,\dots,T$}
    \State Sample a minibatch $\mathcal{S} \subset \{1,\dots,\mathcal{N}\}$
    \For{$i \in \mathcal{S}$}
        \State Optimize local variational parameter:
        $\bm{\varphi}_{i}^{(t)} = \mathsf{E}^{(t-1)}_{g} \left[ \bm{\eta}_l(\bm{\theta}_g, y_i) \right]$
    \EndFor
    \State Calculate $\rho_t = (t+\tau)^{-\chi}$
    \State Compute the intermediate global parameters
    $\widehat{\bm{h}}^{(t)} = \bm{\lambda} +\mathcal{C}\sum_{i \in \mathcal{S}}\left(\mathsf{E}^{(t)}_{l} \left[ \bm{t}(\theta_{l_i}, y_i) \right],1\right)^\top$
    \State Update global variational parameter 
    $\bm{h}^{(t)} = (1 - \rho_t)\bm{h}^{(t-1)} + \rho_t \widehat{\bm{h}}^{(t)}$
\EndFor
\Statex \textbf{Output} Optimal variational parameters $\bm{h},\bm{\varphi}$
\end{algorithmic}
\label{alg:SVI_minibatch}
\end{algorithm}